Die Luftwiderstandskraft ist (FoTaBe S.~165)
\begin{align}
 F_\text{W} &= \frac{1}{2}\,c_w\,\rho\,A\,v^2
\end{align}
Der Wind trifft die Stange quer, die Stirnfläche ist das Rechteck $A = l\cdot d = \AFaP{0}{2}$. Für einen quer angeströmten Kreiszylinder mit $l/d \approx \ldFP{0}{2}$ gibt die Tabelle der Widerstandszahlen (FoTaBe S.~189, Wert für $l/d = 40$) $c_w = \cwZO$; die Dichte der Luft ist $\rho = \rhLO$ (FoTaBe S.~187). Damit ist
\begin{align}
 F_\text{W} &= \FFaF \\
   &= \frac{1}{2}\cdot\cwZ\cdot\rhL\cdot\lFa\cdot\dFa\cdot\left(\vWi\right)^2 \\
   &= \FFa \approx \result{\FFaP{0}{2}}
\end{align}
